\documentclass[11pt]{article}
\usepackage[margin=1in]{geometry}
\usepackage{hyperref}
\usepackage{enumitem}
\title{Tinkercad Design and Revision}
\author{Piter Garcia}
\date{\today}
\begin{document}
\maketitle

\section*{Grade and Class Match}
\textbf{Grade:} Grade 5\\
\textbf{Likely blocks:} Day 2 · 10:45-11:35 · Pickett; Day 3 · 10:45-11:35 · McGlashon; Day 5 · 10:45-11:35 · Pecora

\section*{Lesson Focus}
3D design through shape placement, grouping, subtraction, and revision language.

\section*{Learning Targets}
\begin{itemize}[leftmargin=*]
\item Students navigate the Tinkercad workspace confidently.
\item Students use grouping and hole tools to modify a design.
\item Students explain one revision they made and why.
\end{itemize}

\section*{Materials}
\begin{itemize}[leftmargin=*]
\item Student devices
\item Tinkercad accounts
\item Projected example model
\end{itemize}

\section*{Suggested Lesson Flow}
\begin{enumerate}[leftmargin=*]
\item Open with the exact device, tool, or routine students need in order to start successfully.
\item Model one example task before students begin independent or partner work.
\item Give students time to build, code, design, or document their work while circulating for support.
\item Close with a short share-out, reflection, or debugging discussion.
\end{enumerate}

\section*{Source Notes}
This lesson packet is enriched with transcript-derived lesson notes, the school schedule roster, and official starting-point resources. See the paired Markdown guide for clickable links.

\bibliographystyle{plain}
\bibliography{tinkercad-design-and-revision}
\end{document}
